\phantomsection
\chapter*{Melodi dan Kotak Musik Tua}
\addcontentsline{toc}{section}{Melodi dan Kotak Musik Tua --- Ardanis Nabil Z. D.}
\penulis{Ardanis Nabil Z. D.}

\awalcerpen{G}{erimis tipis mulai menetes}, membasahi kaca jendela kafe tua bernuansa kayu di sudut kota. Suasana di dalam ruangan terasa hangat dan tenang, berbanding terbalik dengan rusuhnya jalanan di luar. Bau harum kopi hangat bercampur dengan aroma kue cokelat menciptakan suasana yang nyaman. Di salah satu sudut kafe yang paling redup, Nabil dan Lala duduk bersampingan.

Di meja kayu di antara mereka, terletak sebuah kotak musik kecil berukiran daun semanggi yang tampak anggun. Ketika Nabil memutar kunci kecil di bagian sampingnya, alunan denting lagu \(\textit{Happy Birthday to You}\) berdenting lembut. Nadanya jernih, mengalun amat rapi melintasi waktu, membawa ingatan Nabil kembali jauh ke masa lalu—ke awal mula dari segala sesuatu yang mereka bangun bersama.

Hubungan mereka tidak dimulai dengan perayaan kembang api atau drama romantis yang megah. Semuanya berawal dari sebuah ruangan sekretariat organisasi sekolah yang berantakan. Saat itu, Nabil adalah anggota divisi keamanan yang cenderung tenang, lebih suka bekerja di lapangan tanpa banyak bersantai. Di sisi lain, Lala adalah anggota baru yang energetik, cerdas, dan selalu membawa keceriaan di mana pun ia berada.

Pertemuan kecil itu terjadi di suatu sore yang melelahkan menjelang acara besar sekolah. Semua orang sibuk dengan urusannya masing-masing, sementara Nabil terduduk lelah di sudut ruangan sambil membereskan tumpukan kabel. Tiba-tiba, sebuah gelas plastik berisi teh manis dingin diberikan kepadanya.

``Minum dulu, nih. Muka kamu udah kusam banget kayak berkas-berkas tua di pojokan itu,'' ucap sebuah suara tenang.

Nabil mendongak dan melihat Lala berdiri di sana, tersenyum lebar dengan mata berbinar. Itulah pertama kalinya Nabil benar-benar memperhatikan Lala. Sejak hari itu, interaksi kecil mulai terjalin secara alami. Lala yang selalu kebingungan membawa barang berat sering meminta bantuan Nabil, sementara Nabil yang selalu lupa makan saat bekerja selalu mendapat teguran halus—lengkap dengan camilan—dari Lala.

Ada satu hal tentang Lala yang perlahan tapi pasti berhasil mencuri hati Nabil: suaranya. Lala memiliki kebiasaan bernyanyi kecil saat sedang fokus mengerjakannya sesuatu. Suaranya tidak seperti penyanyi profesional di televisi, tetapi memiliki karakter yang amat lembut, hangat, dan menenangkan. Suara itu mengingatkan Nabil pada lagu-lagu Nadin Amizah—penuh kelembutan, puitis, dan seolah membawa kehangatan di tengah dinginnya malam. Setiap kali Lala bersenandung pelan menyanyikan alunan melodi saat merapikan dokumen, Nabil selalu menemukan dirinya terpaku, diam-diam mendengarkan dan mengagumi keindahan yang terpancar dari sosok di depannya.

Kehadiran Lala yang semula hanyalah sekadar rekan satu organisasi perlahan berubah menjadi alasan utama Nabil untuk selalu hadir lebih awal. Rasa tertarik yang sederhana itu tumbuh secara perlahan menjadi benih-benih cinta tanpa perlu dipaksa.

Ketika akhirnya mereka memutuskan untuk berjalan bersama sebagai sepasang kekasih, kisah cinta yang terjalin sangatlah sederhana namun bermakna. Tidak ada keharusan untuk memberi kejutan mewah setiap minggu, tidak ada tuntutan untuk tampil sempurna di depan publik. Hubungan mereka mengalun bagaikan sungai yang tenang di sore hari.

Nabil dan Lala belajar untuk saling melengkapi dalam hal-hal kecil. Mereka menemukan kebahagiaan dari momen-momen sederhana: menghabiskan malam minggu dengan menyusuri jalanan kota di atas motor tua Nabil, makan bakso di pinggir jalan sambil menertawakan hal-hal konyol, atau sekadar duduk berdua di perpustakaan tanpa bersuara, hanya menikmati kehadiran satu sama lain.

Lala adalah tipe wanita yang praktis dan perhatian, sementara Nabil adalah sosok pengamat yang sabar. Lala mengajari Nabil untuk lebih terbuka dalam mengekspresikan perasaan, sedangkan Nabil memberikan ruang aman bagi Lala ketika wanita itu merasa lelah dengan dunia luar. Seperti alunan lagu favorit Lala yang sering ia senandungkan saat hujan turun, cinta mereka terasa begitu alami, tidak terburu-buru, dan menyembuhkan.

Masa-masa terindah dalam hubungan mereka bukan diukur dari seberapa sering mereka pergi berlibur ke tempat mahal, melainkan dari bagaimana mereka saling menjaga dalam kesederhanaan. Nabil selalu ingat bagaimana Lala akan menyanyikan lagu-lagu berirama tenang saat ia sedang stres menjalani pekerjaan sekolah, dan bagaimana senyum Lala selalu mampu menghapus seluruh rasa lelahnya dalam sekejap.

Tahun-tahun berlalu dengan cepat hingga mereka memasuki tahun ketiga bersama. Di hari ulang tahun Nabil yang ke-17, Lala sudah menyiapkan sebuah kejutan khusus. Hari itu, hujan turun deras sejak pagi, memadamkan aliran listrik di kawasan rumah Nabil.

Di tengah kegelapan yang hanya diterangi sebatang lilin kecil di atas meja kafe, Lala duduk di hadapan Nabil. Wajahnya tersinari cahaya lilin yang memancar, membuatnya terlihat begitu anggun dan berharga di mata Nabil. Dari dalam tasnya, Lala mengeluarkan sebuah kotak kayu tua berukuran genggaman tangan dengan ukiran daun semanggi yang sangat detail di permukaannya.

``Ini buat kamu,'' kata Lala pelan, menyerahkan kotak itu.

Nabil menerima kotak kayu tersebut dengan rasa ingin tahu. Ketika ia memutar kunci kecil di bagian sampingnya, alunan lembut \textit{Happy Birthday to You} berdenting dari dalam kotak. Suaranya jernih, bergema indah di dalam ruangan yang gelap dan dingin.

``Ini kotak musik tua yang aku temukan di pasar antik,'' bisik Lala, menatap mata Nabil dengan penuh kehangatan. ``Aku minta pengrajinnya buat memperbaiki mesin di dalamnya dan mengganti melodi standarnya jadi lagu ini. Aku sengaja pilih aransemen yang sedikit klasik dan tenang.''

Lala kemudian tersenyum manis dan mulai ikut bernyanyi mengikuti denting kotak musik itu. Suaranya yang merdu menyatu dengan harmoni kayu tua tersebut, mengisi keheningan malam yang gelap dan dingin. Bagi Nabil, suara Lala yang mengalun di bawah remang cahaya lilin itu adalah pertunjukan musik paling indah yang pernah ia dengar sepanjang hidupnya. Indahnya kebersamaan mereka malam itu melampaui segala kemewahan dunia.

Kini, bertahun-tahun telah berlalu sejak sore itu. Nabil menatap kotak musik yang masih berdenting lembut di atas meja kafe, lalu mengalihkan pandangannya pada wanita yang duduk di hadapannya. Lala masih sama seperti dulu—masih dengan senyum hangat yang sama, mata berbinar yang sama, dan keindahan yang tak pernah memudar oleh waktu.

Nabil mengulurkan tangannya, menggenggam jemari Lala yang terasa hangat di atas meja.

``Kamu tahu, La?'' suara Nabil terdengar dalam dan tulus, memecah keheningan di antara mereka. ``Setiap kali aku mendengarkan kotak musik ini, aku nggak cuma ingat ulang tahunku. Aku selalu diingatkan tentang seberapa beruntungnya aku memiliki kamu dalam hidupku.''

Lala menaikkan alisnya tipis, tersenyum kecil menanggapi ucapan Nabil. ``Tiba-tiba banget jadi romantis gini?''

``Eh serius buset dah,'' lanjut Nabil, matanya menatap hangat. ``Di dunia yang berjalan begitu cepat dan sering bikin capek ini, kamu adalah tempat aku pulang. Aku selalu bersyukur sama Tuhan karena udah mempertemukan kita di sekolah yang berantakan itu.''

Nabil mengusap lembut punggung tangan Lala. ``Melihat kamu selama ini, mendengar suara kamu bernyanyi saat hujan, dan melihat bagaimana kamu selalu memperlakukan hal-hal kecil dengan penuh cinta$\dots$ itu semua adalah keindahan yang nggak pernah berhenti bikin aku kagum. Kamu adalah hal terbaik yang pernah hadir dalam hidupku, La. Rasa syukurku karena memiliki kamu nggak akan pernah cukup cuma diungkapkan dengan kata-kata.''

Pipi Lala memerah mendengar ketulusan dalam ucapan Nabil. Matanya menipis manis, memancarkan rasa bahagia yang meluap dari dalam dirinya.

Lala membalas ``Aku juga bersyukur, Bil,'' bisik Lala pelan. ``Terima kasih sudah selalu jadi pendengar yang sabar untuk semua lagu dan cerita hidupku.''

Denting kotak musik tua itu akhirnya perlahan usai, meninggalkan keheningan yang amat nyaman di antara mereka berdua. Di luar, hujan gerimis masih setia membasahi halaman, namun di dalam kafe kecil itu, sepasang kekasih yang saling mengikat telah menemukan harmoni mereka—sebuah melodi cinta yang tidak akan pernah pudar ditelan waktu.

\jedahias
